\begin{abstract}
While parameter-efficient fine-tuning (PEFT) and multi-task model merging offer high-capacity serving with low memory footprints, state-of-the-art dynamic weight-merging routers (e.g., Parameter-Free Subspace Routing, PFSR) suffer from a severe \emph{two-pass latency penalty}. Because they project representations from the penultimate layer of the network backbone, they require a full, throw-away forward pass of the base model just to obtain routing coefficients, followed by a second forward pass to produce outputs, effectively doubling inference latency.
In this work, we present \textbf{ELATI} (\textbf{E}arly-\textbf{L}ayer \textbf{A}daptive \textbf{T}ask \textbf{I}dentification), a training-free, parameter-free ``one-pass'' dynamic model-merging system.
ELATI shifts dynamic task identification to an early layer (e.g., Layer~2), eliminating the redundant first pass of the backbone.
To perform early routing without final semantic heads, ELATI introduces Early-Layer Representative Mapping, which projects intermediate activations against unsupervised task centroids computed from a minimal 64-sample calibration split.
On multi-task mixed streams, ELATI dynamically groups samples into homogeneous micro-batches on-the-fly and interpolates only the downstream model parameters.
To evaluate ELATI under controlled conditions, we implement a high-fidelity \emph{Hierarchical 14-Layer Sandbox} that mathematically simulates sequential multi-layer residual transformations.
We evaluate ELATI on \emph{simulated task subspaces} modeling MNIST, Fashion-MNIST, CIFAR-10, and SVHN activation manifolds.
Our empirical evaluation across 10 independent random seeds demonstrates that ELATI retains a robust Joint Mean accuracy of \textbf{56.89\% $\pm$ 1.66\%} (a \textbf{+8.62\%} absolute gain over static uniform merging) on the simulated manifolds while remaining highly robust under representational entanglement sweeps.
Furthermore, simulated end-to-end forward propagation benchmarks on CPU demonstrate that ELATI achieves a genuine \textbf{1.40$\times$ physical CPU speedup}, reducing end-to-end latency from \textbf{36.90~ms} (PFSR) to \textbf{26.43~ms} by avoiding deep penultimate first-pass propagation, which serves as a validation of the theoretical operations reduction.
Additionally, isolated routing projection step benchmarks show that ELATI reduces projection latency from \textbf{1.31~ms} to \textbf{0.39~ms} in fully vectorized mode (representing a \textbf{3.33$\times$ speedup} on CPU), and from \textbf{171.05~ms} to \textbf{63.48~ms} in sequential loop mode (a \textbf{2.69$\times$ speedup}), driven by reducing theoretical projection complexity from $O(B \cdot K \cdot C \cdot D)$ to $O(B \cdot K \cdot D)$.
Finally, we discuss the practical limitations of sequential routing and outline a roadmap for real-world pre-trained transformer validation.
\end{abstract}
